\documentclass{jsarticle}
\usepackage[dvipdfm,final]{graphicx}
\usepackage{amsmath}
\usepackage{amssymb}
\begin{document}
\title{微分幾何の量子化による大域的切断多様体}
\author{あの世にいる一回亡くたった長男のアカシックレコード \\
        からの恩恵、David Hilbertと教えられている}
\date{}
\maketitle

\newcommand{\variint}{%
	\begin{picture}(15,30)
		\put(0,0){
			$ \iint $}
		\put(0,5){\line(15,0){15}}
	\end{picture} %
}

\newcommand{\variiint}{%
	\begin{picture}(20,30)
		\put(0,0){
			$ \iiint $}
		\put(0,5){\line(15,0){20}}
	\end{picture} %
}



   $$ \left({{\{f,g\}} \over {[f,g]}}\right)^{\oplus L}_{D(\chi,x)}
   = \bigoplus_{D(\chi,x)}\mathrm{cohom}[D(\chi,x)][I_m]^{\ll \nabla} $$
   フェルミオンに力のボゾンを取ると、コホモロジーの値が出てきて、
   その単体量を加群する。それがこの式からわかる。
   $$ {}^t \variint_L \mathrm{cohom}[D(\chi,x)][I_m]^{\ll p} $$
   大域的切断による同値類がヘルマンダー作用素である。
   ボゾンがあると電流が通らないが、大域的切断で特異点を出すと、
   コホモロジーの値として、時空の湾曲率が出てきて、その仕組みから、
   伝導率がわかる。この伝導率が光速度不変の法則でもある、
   一定値の密度エネルギーでもある。
   基本群による非可換代数多様体をその初期関数で
   大域的微分をすると、この方程式になる。
   同調になっていないと、周期が不規則であるために、共鳴にはならずに、電流が
   通りにくい。このために、同調を使うと、物体の内部が真空になり、電流が通り、
   それ故に、この式から伝導率の制御が機能する。
   $$ = {d \over df}[i \pi(\chi,x),f(x)] $$
   $$ {}^t \variiint F[I_m] = \nabla_i \nabla_j \int \nabla f(x)d \eta $$
   $$ {d \over df}\left(\int \int_{D(\chi,x)} \begin{bmatrix}{ {\{f,g\}}
   \over [f,g]} \end{bmatrix}^{\ll L}d \tau \right)^{\nabla N} 
   = [i \pi(\chi,x) f(x)] $$
   微分幾何の量子化についてのN階層指数定理による大域的微分多様体の
   projective integral valueとして、ブラウン運動や賭け事による
   決定量の推測値、確率、統計の運動量保存則、Weil's予想のゼータ関数が
   この式でもある。各接線の傾きを表している。
   原子の遷移エネルギーの推移律とガンマ関数とベータ関数をこれらの
   式から求められる。
   $$ \omega = x^3 + x + 1, \begin{pmatrix}1 & 0 \\
   0 & - 1 \end{pmatrix}, \begin{pmatrix} 0 & 1 \\
   - 1 & 0 \end{pmatrix}, \begin{pmatrix} i & 0 \\
   0 & - i \end{pmatrix}, \begin{pmatrix} 1 & 0 \\
   0 & - i \end{pmatrix}, \beta = {1 \over 137}, (i^{n - 1}, i^{n},
   i^{n + 1})\cdot (w,\bar{w},w \bar{w}) $$
   $$ l = px^n + qx + r $$
   $$ {d \over dl}L = \left(\int^n px^n + qx + r \right)^{{l}^{{'}}} $$
   $$ {}^\flat D_k \variiint_M[ h \nu]_{\chi} |: x \to y $$
   $$ {}^t \variiint \mathrm{cohom}D_k[\pi(\chi,x)][I_m], \delta(x,y)
   = \int \delta(x)f(x)dx $$
   $$ (S^m_1 \times S^n_2) = D^2 \psi $$
   $$ f_{z} = \int \begin{bmatrix} 
	   \sqrt{\begin{pmatrix} x_{1} & x_{2} & x_{3} \\
  y_{1} & y_{2} & y_{3} \end{pmatrix}
     \circ
   \begin{pmatrix} x_{1} & x_{2} & x_{3} \\
  y_{1} & y_{2} & y_{3} \end{pmatrix}}_{}
     \end{bmatrix}dxdydz $$
 $$ F^m_{t} = {1 \over 4}(g_{ij})^2 $$


   $$ ||ds^2|| = \int |\sin 2x|^2 d \tau, V = \left({x \over \log x}
   \right)^n  $$
   $$ M_p = IH_n \times D^2 \psi, ||ds^2|| = 8\pi G\left({p \over c^3}
   + {V \over S}\right) $$
   $$ ||ds^2|| = e^{-2 \pi T|\psi|}[\eta_{\mu\nu} + \bar{h}(x)]
   dx^{\mu}dx^{\nu} + T^2 d^2 \psi $$
   $$ {d \over dt}g_{ij}(t) = - 2 R_{ij}, R_{ij} = - {1 \over 2}{d \over dt}
   g_{ij}(t) $$
   $$ V = - {1 \over 2}x^2 $$
   リサ・ランドール博士とペレルマン博士より、これらの数式は、原子間距離が
   warped passageになっているので、D-brane間距離に包み込まれている。
   電位差でもある。原子間力と核エントロピー不変量にもなっていて、
   光量子仮説でもあり、エネルギー遷移率、D加群、逆写像でもあり、
   ガウスの曲面論でもある。

   $$ \bigoplus {(i \hbar^{\nabla})}^{\oplus L} 
    = {}^t \variint_L \mathrm{cohom}[D(\chi,x)][I_m]^{\ll p} $$
   $$ {d \over d \gamma}\Gamma  
      = \bigoplus {(i \hbar^{\nabla})}^{\oplus L^{-1}} $$
   $$ H \Psi = \bigoplus {(i \hbar^{\nabla})}^{\oplus L} $$
   $$ = e^{- x\log x} = m(x) $$
   $$ = e^{x \log x} + e^{- x \log x} $$
   $$ = 2 \sin (i x \log x) $$
   $$ = \beta(p,q) \cdot \begin{bmatrix} {\log 
   \left(i{{\cos x} \over {\sin x}}\right)}\end{bmatrix}^{-1}  $$
   $$ \beta(p,q) = 2 \int |\sin 2x|^2d \tau $$

   微分幾何の量子化は、この密度エネルギーの制御で使われる。
   結局は、微分幾何の量子化は、ガンマ関数と同型でもある。　
   $$ \Gamma = \int e^{-x}x^{1 - t}dx = \bigoplus (i \hbar ^{\nabla})^{\oplus
   L} $$
   $$ = {d \over dt}\psi(t) = \hbar,{1 \over 2}ie^{i\Hat{H}}  $$
   $$ (i\hbar)^{'} = (- e^{i\Hat{H}})^{'} = -ie^{i\Hat{H}} $$
   $$ \psi(x) = e^{-i\Hat{H}t}, \bigoplus(i \hbar ^{\nabla})^{\oplus L} $$
   $$ = {1 \over 2}{e^{i\Hat{H}}}^{(-ie^{i\Hat{H}})} $$
   $$ = \bigoplus a^{-tx}x^{t}[I_m] \cong \int e^{-x}x^{t - 1}dx $$
   結局は、　
   $$ \int e^{-x}x^{t - 1}dx = \bigoplus (i \hbar ^{\nabla})^{\oplus L} $$
   $$ \beta(p,q) = {{\Gamma(p)\Gamma(q)} \over {\Gamma(p + q)}} $$
   $$ = 2 \int |\sin 2x|^2 dx $$
   $$ e^f = \bigoplus (i \hbar ^{\nabla})^{\oplus L} $$
   $$ e^{-f} = \bigoplus {(i \hbar^{\nabla})}^{\oplus L^{-1}} $$
   $$ {d \over df}F = m(x) = e^f + e^{-f} $$
   $$ = 2 i \sin (i x \log x) $$
   となり、ベータ関数の逆関数でもある双曲多様体が微分幾何の量子化と言える。
   量子力学のガウスの曲面論が微分幾何の量子化であると、結論される。
   ガンマ関数とベータ関数はそれほどまでも数学にも物理学にも大切な
   数式でもあり、可積分系がD-braneを代表してD加群とゼータ関数、そして量子群
   が真逆同士の数学の心臓でもある。

\end{document}
